\documentclass[11pt,a4paper]{article}

\usepackage{fontspec}
\setmainfont{PTF55F.ttf}[BoldFont=PTF75F.ttf,ItalicFont=PTF56F.ttf,BoldItalicFont=PTF76F.ttf]
\setsansfont{PTS55F.ttf}[BoldFont=PTS75F.ttf,ItalicFont=PTS56F.ttf,BoldItalicFont=PTS76F.ttf]
\setmonofont{PTM55F.ttf}[Scale=0.92]
\usepackage[russian]{babel}
\AtBeginDocument{\shorthandoff{"}}

\usepackage[margin=2.2cm]{geometry}
\usepackage{xcolor}
\usepackage{minted}
\usepackage[most,minted]{tcolorbox}
\usepackage{booktabs}
\usepackage{tabularx}
\usepackage{enumitem}
\usepackage{titlesec}
\usepackage[colorlinks=true,linkcolor=blue!50!black,urlcolor=blue!50!black]{hyperref}

\setlength{\parindent}{0pt}
\setlength{\emergencystretch}{3em}
\setlength{\parskip}{0.55em}
\setlist{itemsep=0.15em,topsep=0.3em}
\titleformat{\section}{\Large\sffamily\bfseries}{\thesection}{0.6em}{}
\titleformat{\subsection}{\large\sffamily\bfseries}{\thesubsection}{0.6em}{}
\titleformat{\subsubsection}{\normalsize\sffamily\bfseries}{\thesubsubsection}{0.6em}{}

\definecolor{codebg}{HTML}{F6F7F9}

% Листинги: minted (нужен pygmentize и tectonic -Z shell-escape).
\usemintedstyle{friendly}
% Пробелы в листингах выводятся настоящим глифом U+0020, чтобы отступы
% сохранялись при копировании кода из PDF.
\newcommand{\copyspace}{\XeTeXglyph\XeTeXcharglyph"0020\relax}
\tcbset{codebox/.style={listing only, breakable, enhanced, colback=codebg, frame hidden,
  arc=1pt, left=4pt, right=4pt, top=1pt, bottom=1pt, before skip=0.5em, after skip=0.6em}}
\newtcblisting{py}{codebox, listing engine=minted, minted language=python, minted options={fontsize=\small,breaklines,tabsize=4,showspaces,space=\copyspace}}
\newtcblisting{shell}{codebox, listing engine=minted, minted language=console, minted options={fontsize=\small,breaklines,tabsize=4,showspaces,space=\copyspace}}
\newtcblisting{toml}{codebox, listing engine=minted, minted language=toml, minted options={fontsize=\small,breaklines,tabsize=4,showspaces,space=\copyspace}}
\newtcblisting{plaintext}{codebox, listing engine=minted, minted language=text, minted options={fontsize=\small,breaklines,tabsize=4,showspaces,space=\copyspace}}
\newtcblisting{plainwide}{codebox, listing engine=minted, minted language=text, minted options={fontsize=\scriptsize,breaklines,tabsize=4,showspaces,space=\copyspace}}
\newcommand{\code}[1]{\texttt{#1}}

\begin{document}

\begin{center}
  {\sffamily\small Программирование на языке Python · 2 курс}\\[0.6em]
  {\sffamily\LARGE\bfseries Лабораторная работа №2}\\[0.4em]
  {\sffamily\Large ООП и производительность.\\ Задача «Миллиард строк» (продолжение)}
\end{center}

\tableofcontents
\newpage

%=====================================================================
\section{Общие положения}

\subsection{Цель работы}

Работа продолжает лабораторную работу №1. Цель работы~— изучение объектно-ориентированного
программирования в Python и методов измерения производительности. По итогам работы студент
должен:

\begin{itemize}
  \item уметь проектировать классы с использованием магических методов, свойств и \code{dataclass};
  \item уметь измерять время выполнения программы и находить узкие места с помощью профилировщика;
  \item понимать, от чего зависит производительность программы на Python.
\end{itemize}

%=====================================================================
\section{Объектно-ориентированное программирование}

\subsection{Классы и объекты}

\begin{py}
class Point:
    dimensions = 2                      # атрибут класса (аналог static)

    def __init__(self, x: float, y: float) -> None:
        self.x = x                      # атрибуты экземпляра создаются присваиванием
        self.y = y

    def norm(self) -> float:
        return (self.x ** 2 + self.y ** 2) ** 0.5


p = Point(3, 4)         # оператор new не используется
print(p.norm())         # 5.0
print(Point.dimensions, p.dimensions)
\end{py}

\begin{itemize}
  \item \code{\_\_init\_\_}~— инициализатор. Перегрузка конструкторов
        не поддерживается; вместо неё используются значения аргументов по умолчанию или
        альтернативные конструкторы на основе \code{@classmethod}.
  \item \code{self}~— явный аналог \code{this}. Он передаётся первым аргументом в каждый метод;
        обращение к атрибутам выполняется только через него: \code{self.x}, а не \code{x}.
  \item Поля в теле класса не объявляются: атрибут появляется при первом присваивании.
        Принято создавать все атрибуты в \code{\_\_init\_\_}.
  \item Деструктор в привычном для C++ смысле отсутствует; освобождение ресурсов выполняется
        с помощью \code{with}.
\end{itemize}

\subsection{Инкапсуляция}

Модификаторы доступа в Python отсутствуют; вместо них используются соглашения об именах:
\begin{itemize}
  \item \code{self.name}~— публичный атрибут (public);
  \item \code{self.\_name}~— внутренний атрибут (protected), не предназначенный для использования вне класса.
        Технически он доступен, но обращение к нему извне считается нарушением соглашения;
  \item \code{self.\_\_name}~— приватный атрибут (private); использование из вне через прямое имя
        невозможно~— имя атрибута подвергается искажению (name mangling) в \code{\_Class\_\_name}
        для предотвращения конфликтов имён в подклассах. Применяется редко.
\end{itemize}

Вместо пар методов \code{getX()/setX()} используются \textbf{свойства} (properties): для внешнего кода
свойство выглядит как обычный атрибут, но при обращении к нему вызывается метод.

\begin{py}
class Temperature:
    def __init__(self, celsius: float) -> None:
        self.celsius = celsius          # вызывается сеттер, определённый ниже

    @property
    def celsius(self) -> float:
        return self._celsius

    @celsius.setter
    def celsius(self, value: float) -> None:
        if value < -273.15:
            raise ValueError("температура ниже абсолютного нуля")
        self._celsius = value

    @property
    def fahrenheit(self) -> float:      # свойство только для чтения
        return self._celsius * 9 / 5 + 32


t = Temperature(20)
t.celsius = 25          # проверка выполняется в сеттере
print(t.fahrenheit)     # 77.0
# t.fahrenheit = 0      # AttributeError: сеттер не определён
\end{py}

\subsection{Магические методы}

Методы, имена которых начинаются и заканчиваются двойным подчёркиванием (\emph{dunder}-методы),
вызываются интерпретатором при участии объекта во встроенных операциях. Механизм аналогичен
перегрузке операторов в C++, но охватывает более широкий круг операций: преобразование в строку,
получение длины, итерацию, проверку вхождения и~т.\,д.

\begin{tabularx}{\linewidth}{@{}>{\ttfamily}l>{\ttfamily}lX@{}}
\toprule
\textrm{\textbf{Метод}} & \textrm{\textbf{Вызывается при}} & \textbf{Аналог в C++} \\
\midrule
\_\_repr\_\_(self) & repr(x), в REPL, в отладчике & --- \\
\_\_str\_\_(self) & str(x), print(x), f"\{x\}" & \code{operator<<} \\
\_\_eq\_\_(self, other) & x == y & \code{operator==} \\
\_\_lt\_\_(self, other) & x < y, sorted() & \code{operator<} \\
\_\_hash\_\_(self) & ключ dict, элемент set & \code{std::hash} \\
\_\_len\_\_(self) & len(x) & \code{size()} \\
\_\_getitem\_\_(self, i) & x[i] & \code{operator[]} \\
\_\_setitem\_\_(self, i, v) & x[i] = v & \code{operator[]} (lvalue) \\
\_\_contains\_\_(self, v) & v in x & --- \\
\_\_iter\_\_(self) & for v in x & \code{begin()/end()} \\
\_\_add\_\_(self, other) & x + y & \code{operator+} \\
\_\_iadd\_\_(self, other) & x += y & \code{operator+=} \\
\_\_mul\_\_(self, other) & x * y & \code{operator*} \\
\_\_bool\_\_(self) & if x: & \code{explicit operator bool} \\
\bottomrule
\end{tabularx}

\medskip
Пример~— вектор на плоскости:

\begin{py}
class Vec2:
    def __init__(self, x: float = 0.0, y: float = 0.0) -> None:
        self.x, self.y = x, y

    def __repr__(self) -> str:
        return f"Vec2({self.x!r}, {self.y!r})"

    def __eq__(self, other: object) -> bool:
        if not isinstance(other, Vec2):
            return NotImplemented       # сравнение с этим типом не поддерживается
        return self.x == other.x and self.y == other.y

    def __add__(self, other: "Vec2") -> "Vec2":
        return Vec2(self.x + other.x, self.y + other.y)

    def __mul__(self, k: float) -> "Vec2":      # v * 3
        return Vec2(self.x * k, self.y * k)

    def __rmul__(self, k: float) -> "Vec2":     # 3 * v
        return self * k

    def __abs__(self) -> float:
        return (self.x ** 2 + self.y ** 2) ** 0.5

    def __iter__(self):                         # позволяет распаковку x, y = v
        yield self.x
        yield self.y


v = Vec2(1, 2) + 3 * Vec2(1, 0)
print(v)                # Vec2(4, 2)
x, y = v
\end{py}

\begin{itemize}
  \item \code{\_\_repr\_\_} возвращает однозначное представление объекта для разработчика,
        в идеале~— выражение, по которому объект можно воссоздать. Метод рекомендуется
        определять во всех классах.
  \item Если \code{\_\_str\_\_} не определён, функция \code{print} использует \code{\_\_repr\_\_}.
  \item При определении \code{\_\_eq\_\_} хеш-функция по умолчанию отключается, и объект
        не может быть ключом словаря, пока не определён \code{\_\_hash\_\_}.
  \item Ключевое слово \code{yield} превращает функцию в \emph{генератор}, выдающий значения
        по одному по мере запроса. Это наиболее простой способ реализовать \code{\_\_iter\_\_}.
\end{itemize}

\subsection{Наследование}

\begin{py}
class Shape:
    def __init__(self, name: str) -> None:
        self.name = name

    def area(self) -> float:
        raise NotImplementedError       # аналог чисто виртуального метода

    def __repr__(self) -> str:
        return f"{type(self).__name__}({self.name!r}, area={self.area():.2f})"


class Circle(Shape):
    def __init__(self, r: float) -> None:
        super().__init__("круг")        # вызов инициализатора базового класса
        self.r = r

    def area(self) -> float:            # переопределение метода
        return 3.14159 * self.r ** 2


class Square(Shape):
    def __init__(self, a: float) -> None:
        super().__init__("квадрат")
        self.a = a

    def area(self) -> float:
        return self.a ** 2


shapes = [Circle(1), Square(2)]
print(sum(s.area() for s in shapes))
print(isinstance(shapes[0], Shape))    # True
\end{py}

Все методы в Python ведут себя как виртуальные: вызов \code{s.area()} всегда разрешается
по фактическому классу объекта, ключевые слова \code{virtual} и \code{override} отсутствуют.
Поддерживается множественное наследование; порядок поиска методов (MRO) доступен через
атрибут \code{Circle.\_\_mro\_\_}.

\textbf{Утиная типизация.} Для использования объекта в функции не требуется его наследование
от определённого класса~— достаточно наличия необходимых методов. Объект с методами
\code{\_\_len\_\_} и \code{\_\_getitem\_\_} поддерживает \code{len()} и индексацию, объект
с методом \code{\_\_iter\_\_}~— обход циклом \code{for}.

\subsection{Методы класса и статические методы}

\begin{py}
class Date:
    def __init__(self, year: int, month: int, day: int) -> None:
        self.year, self.month, self.day = year, month, day

    @classmethod
    def from_string(cls, s: str) -> "Date":     # альтернативный конструктор
        year, month, day = (int(part) for part in s.split("-"))
        return cls(year, month, day)

    @staticmethod
    def is_leap(year: int) -> bool:             # функция в пространстве имён класса
        return year % 4 == 0 and (year % 100 != 0 or year % 400 == 0)


d = Date.from_string("2026-09-01")
Date.is_leap(2028)
\end{py}

\subsection{Классы данных и \texttt{\_\_slots\_\_}}

Для классов, предназначенных преимущественно для хранения данных, декоратор \code{@dataclass}
генерирует методы \code{\_\_init\_\_}, \code{\_\_repr\_\_} и \code{\_\_eq\_\_} по аннотациям полей:

\begin{py}
from dataclasses import dataclass, field

@dataclass
class Student:
    name: str
    group: int
    grades: list[int] = field(default_factory=list)    # не grades=[]

    def average(self) -> float:
        return sum(self.grades) / len(self.grades) if self.grades else 0.0


s = Student("Анна", 2)
print(s)        # Student(name='Анна', group=2, grades=[])

@dataclass(frozen=True)         # неизменяемый класс, получает __hash__
class Coord:
    lat: float
    lon: float
\end{py}

По умолчанию атрибуты объекта хранятся в словаре \code{obj.\_\_dict\_\_}, что позволяет
добавлять их динамически, но увеличивает расход памяти и время доступа. Объявление
\code{\_\_slots\_\_} заменяет словарь фиксированным набором полей, аналогично \code{struct} в C++.
В примере ниже объём памяти измеряется модулем \code{tracemalloc} по 100\,000 объектов
(в значение входит и 8-байтная ссылка на объект в списке); функция \code{sys.getsizeof}
для такого сравнения не подходит, поскольку не учитывает память, занятую атрибутами.

\begin{py}
import tracemalloc
from dataclasses import dataclass


class Pixel:                    # атрибуты хранятся в __dict__
    def __init__(self, x: int, y: int, color: int) -> None:
        self.x, self.y, self.color = x, y, color


class SlotPixel:                # фиксированный набор полей
    __slots__ = ("x", "y", "color")

    def __init__(self, x: int, y: int, color: int) -> None:
        self.x, self.y, self.color = x, y, color


p = SlotPixel(1, 2, 0xFF0000)
# p.alpha = 1    # AttributeError: атрибут отсутствует в __slots__


@dataclass(slots=True)          # аналогичный результат для dataclass
class Pixel2:
    x: int
    y: int
    color: int


def bytes_per_object(cls: type, n: int = 100_000) -> float:
    tracemalloc.start()         # учёт всей памяти, выделяемой интерпретатором
    objects = [cls(1, 2, 3) for _ in range(n)]
    used, _ = tracemalloc.get_traced_memory()
    tracemalloc.stop()
    return used / len(objects)


print(f"Pixel:     {bytes_per_object(Pixel):.0f} байт")
print(f"SlotPixel: {bytes_per_object(SlotPixel):.0f} байт")
# Вывод (Python 3.13, 64-битная платформа):
# Pixel:     104 байт
# SlotPixel: 64 байт
\end{py}

%=====================================================================
\section{Измерение производительности}

\subsection{Время выполнения программы}

\begin{shell}
$ time uv run main.py measurements_10m.txt > /dev/null
real    0m4.512s      # время выполнения программы
\end{shell}

Измерение выполняется несколько раз: первый запуск может быть медленнее, пока файл не загружен
в кеш операционной системы. Время выполнения отдельного участка кода измеряется функцией
\code{time.perf\_counter()}: она возвращает показания монотонного таймера высокого разрешения
в секундах, поэтому значение имеет только разность двух показаний.

\begin{py}
import time

start = time.perf_counter()
total = sum(i * i for i in range(10_000_000))
elapsed = time.perf_counter() - start
print(f"Сумма вычислена за {elapsed:.3f} с")
# Сумма вычислена за 0.494 с
\end{py}

\subsection{Профилирование}

Для определения участков программы, на которые приходится основное время выполнения,
используется профилировщик из стандартной библиотеки:
\begin{shell}
$ uv run python -m cProfile -s tottime main.py measurements_1m.txt | head -20
\end{shell}
Столбец \code{tottime} содержит время выполнения внутри самой функции, \code{ncalls}~— число
вызовов.

Столбец \code{cumtime} содержит время выполнения функции вместе со всеми вызванными из неё
функциями. Профилировщик сам замедляет программу, поэтому его результаты используются для поиска
узких мест, а не для измерения итогового времени.

\subsection{Построчное профилирование}

Профилировщик \code{cProfile} указывает функцию, на которую приходится основное время, но
не строку внутри неё. Время выполнения каждой строки функции измеряет пакет \code{line\_profiler}.
Пакет не входит в стандартную библиотеку и добавляется как зависимость для разработки:
\begin{shell}
$ uv add --dev line_profiler
\end{shell}

Исследуемая функция помечается декоратором \code{@profile}:
\begin{py}
from line_profiler import profile


@profile
def primes(n: int) -> list[int]:
    result = []
    for k in range(2, n):
        is_prime = True
        for d in range(2, int(k**0.5) + 1):
            if k % d == 0:
                is_prime = False
                break
        if is_prime:
            result.append(k)
    return result


if __name__ == "__main__":
    print(len(primes(100_000)))
\end{py}

Без переменной окружения \code{LINE\_PROFILE} декоратор ничего не делает, и программа
выполняется как обычно. При запуске с \code{LINE\_PROFILE=1} результаты измерений записываются
в файл \code{profile\_output.txt} в текущем каталоге:
\begin{shell}
$ LINE_PROFILE=1 uv run primes.py
9592
...
Wrote profile results to profile_output.txt
\end{shell}

Содержимое \code{profile\_output.txt}:
\begin{plainwide}
Timer unit: 1e-09 s

Total time: 0.918304 s
File: /home/student/lab1/primes.py
Function: primes at line 4

Line #      Hits         Time  Per Hit   % Time  Line Contents
==============================================================
     4                                           @profile
     5                                           def primes(n: int) -> list[int]:
     6         1          0.0      0.0      0.0      result = []
     7     99999   17678000.0    176.8      1.9      for k in range(2, n):
     8     99998   14904000.0    149.0      1.6          is_prime = True
     9   2755285  419914000.0    152.4     45.7          for d in range(2, int(k**0.5) + 1):
    10   2745693  421922000.0    153.7     45.9              if k % d == 0:
    11     90406   13286000.0    147.0      1.4                  is_prime = False
    12     90406   13504000.0    149.4      1.5                  break
    13     99998   15337000.0    153.4      1.7          if is_prime:
    14      9592    1755000.0    183.0      0.2              result.append(k)
    15         1       4000.0   4000.0      0.0      return result
\end{plainwide}

Столбец \code{Hits} содержит число выполнений строки, \code{Time}~— суммарное время её
выполнения в единицах \code{Timer unit}, \code{Per Hit}~— среднее время одного выполнения,
\code{\% Time}~— долю от общего времени функции. В приведённом примере более 90\,\% времени
приходится на внутренний цикл (строки~9 и~10). Как и \code{cProfile}, построчный профилировщик
замедляет программу; декораторы \code{@profile} удаляются из кода после завершения анализа.

%=====================================================================
\section{Задание}

Работа продолжает лабораторную работу №1 и выполняется в том же репозитории. Условие задачи
<<Миллиард строк>>, форматы входных и выходных данных и материалы каталога \code{data/} приведены
в описании лабораторной работы №1. Работа состоит из двух частей, обе обязательны.

\subsection{Часть 1. Объектно-ориентированная организация решения}

\begin{enumerate}
  \item \textbf{Класс \code{StationStats}} (файл \code{stats.py}) хранит статистику одной станции
        и предоставляет:
  \begin{itemize}
    \item метод \code{update(value)}~— учёт очередного измерения;
    \item метод \code{merge(other)}~— объединение со статистикой другого объекта \code{StationStats}
          по той же станции;
    \item свойства \code{min}, \code{max}, \code{mean}, \code{count}, доступные только для чтения;
    \item метод \code{\_\_repr\_\_}, возвращающий строку \code{min;mean;max} в формате вывода.
  \end{itemize}
  \item Решение лабораторной работы №1 использует класс \code{StationStats} для хранения
        и объединения статистики станций. Интерфейс командной строки и формат вывода
        не меняются; \code{check.py} по-прежнему выводит \code{OK}.
\end{enumerate}

\subsection{Часть 2. Производительность}

\begin{enumerate}
  \item Для исходной версии решения выполняется профилирование; в
        \code{README.md} приводятся первые строки отчёта и вывод о том, на что приходится
        основное время выполнения.
  \item Решение последовательно ускоряется с контролем времени выполнения. В \code{README.md}
        приводится таблица не менее чем для трёх версий программы: для каждой версии указываются
        внесённые изменения и время обработки файлов из 10~млн и 100~млн строк. Каждая версия
        фиксируется отдельным коммитом.
  \item Все версии остаются корректными (\code{check.py} выводит \code{OK}) и соблюдают требования
        части~1.
\end{enumerate}

%=====================================================================
\section{Порядок сдачи}

До защиты преподавателю направляется ссылка на репозиторий GitHub. Доступность ссылки
проверяется без входа в учётную запись (например, в приватном окне браузера); для приватного
репозитория проверяется наличие преподавателя в списке Collaborators. Сданное состояние
репозитория помечается тегом \code{lab2} и публикуется на GitHub аналогично лабораторной
работе №1; проверке подлежит именно оно.

На защите студент:
\begin{enumerate}
  \item демонстрирует историю коммитов, содержащую отдельные версии решения;
  \item запускает \code{check.py} на файле \code{measurements\_100k};
  \item представляет таблицу замеров и обосновывает влияние каждого изменения на время выполнения;
  \item отвечает на вопросы по коду.
\end{enumerate}

\subsection*{Возможные вопросы}

\begin{enumerate}
  \item Является ли метод \code{\_\_init\_\_} конструктором? Для чего предназначен метод
        \code{\_\_new\_\_}? Каким образом в Python реализуются альтернативные конструкторы?
  \item Каким образом в Python реализуется инкапсуляция? Для чего используется \code{@property}?
  \item Что происходит, если в классе определён \code{\_\_eq\_\_}, но не определён \code{\_\_hash\_\_}?
  \item Чем различаются \code{\_\_repr\_\_} и \code{\_\_str\_\_}? Какой из методов вызывает \code{print}?
  \item Что означает \code{return NotImplemented} в \code{\_\_eq\_\_} или \code{\_\_add\_\_}?
        Для чего нужен \code{\_\_rmul\_\_}?
  \item Чем \code{@classmethod} отличается от \code{@staticmethod} и от обычного метода?
  \item Что такое утиная типизация? Приведите пример из стандартной библиотеки.
  \item Что даёт объявление \code{\_\_slots\_\_} и какие ограничения оно накладывает?
  \item Какое изменение в решении задачи дало наибольший прирост производительности и почему?
\end{enumerate}

\subsection*{Источники}
\begin{itemize}
  \item Модель данных Python (магические методы): \url{https://docs.python.org/3.13/reference/datamodel.html}
  \item Профилировщики Python: \url{https://docs.python.org/3.13/library/profile.html}
  \item Документация line\_profiler: \url{https://kernprof.readthedocs.io/en/latest/manual/index.html}
  \item One Billion Row Challenge: \url{https://github.com/gunnarmorling/1brc}
\end{itemize}

\end{document}
